\chapter{P4: Layanan dan Penyimpanan}
\label{chap:p4}

\section{Rancangan Solusi}

\subsection{Pembagian Domain dan Batas Layanan}

Batas layanan ditarik menurut domain bisnis, bukan menurut lapisan teknis.
Pertanyaan yang dipakai untuk menentukan batas adalah: kemampuan bisnis apa
yang dilayani, data apa yang menjadi miliknya, dan kepada siapa ia bergantung
ketika berjalan. Hasilnya adalah satu domain inti yang menghasilkan data hazard
kanonis dan tiga domain penyampaian yang masing-masing mengonsumsi data itu
dengan semantik berbeda. Identitas dan kontrak publik ditempatkan sebagai
domain pendukung di tepi sistem.

\begin{table}[H]
    \centering
    \small
    \caption{Pembagian domain dan batas layanan}
    \begin{tabularx}{\textwidth}{@{}P{0.20\textwidth}P{0.17\textwidth}X X@{}}
        \toprule
        \textbf{Domain} & \textbf{Layanan} & \textbf{Kemampuan bisnis} & \textbf{Alasan dipisah} \\
        \midrule
        Integrasi sumber (inti) & Aggregator (Go) & Mengambil data BMKG/PVMBG, memetakan ke \code{HazardEvent}, menjadi satu-satunya penulis data kanonis, dan memublikasikan perubahan. & Satu-satunya domain yang memahami format dan kredensial instansi sumber; aturan pemetaan berubah mengikuti sumber, bukan mengikuti klien. \\
        Distribusi real-time & Event Live Feed (Rust) & Mengalirkan event hazard ke klien yang sedang terhubung, dengan filter, proyeksi, dan replay jangka pendek. & Beban berupa banyak koneksi berumur panjang; semantiknya \emph{best-effort} dalam jendela replay, berbeda dari pengiriman yang harus tercatat. \\
        Integrasi pihak ketiga & Webhook Provider (Go) & Mengelola langganan webhook dan mengirim callback bertanda tangan ke sistem eksternal. & Bergantung pada endpoint di luar kendali BNPB yang bisa lambat atau mati; retry dan log pengiriman tidak boleh menahan domain lain. \\
        Peringatan bahaya & Danger Notification (C\#) & Memilih hazard SIAGA/AWAS, menyimpan notifikasi, dan mengirim push melalui FCM. & Memiliki aturan bisnis sendiri (ambang severity, revisi notifikasi) dan dependensi eksternal FCM. \\
        Identitas dan akses (pendukung) & Zitadel, Auth Bridge, sidecar & Login, sesi, penerbitan token, dan pembentukan role/scope. & Menjadi dependensi lintas domain; dipisahkan agar domain bisnis cukup memverifikasi JWT tanpa mengelola kredensial pengguna. \\
        Kontrak publik (pendukung) & Envoy, API Gateway (Go) & Routing, validasi OpenAPI, dan penerjemahan HTTP/SSE ke gRPC. & Bukan pemilik data; memisahkan bentuk API publik dari kontrak internal tiap domain. \\
        \bottomrule
    \end{tabularx}
\end{table}

Tiga domain penyampaian sengaja tidak digabung walaupun ketiganya
mengonsumsi event yang sama. Ketiganya berbeda pada tiga hal: tujuan
penyampaian (klien yang sedang terhubung, sistem eksternal, perangkat
pengguna), jaminan yang diberikan (replay terbatas, pengiriman tercatat dengan
retry, push satu kali per revisi), dan dependensi eksternal (tidak ada,
endpoint callback, FCM). Bila digabung, endpoint webhook yang lambat atau FCM
yang gagal dapat menahan aliran feed. Dengan pemisahan, kegagalan satu
domain hanya menumpuk backlog pada queue miliknya (Bab~\ref{chap:p5}).

Hubungan antardomain hanya melalui dua jenis kontrak: event
\code{HazardEventPublished} dari domain inti ke domain penyampaian, dan gRPC
untuk query, misalnya Danger Notification yang menanyakan status kesegaran
sumber ke Aggregator. Kontrak Protobuf menjadi batas lintas bahasa, sehingga
setiap domain dapat memakai bahasa yang paling sesuai: Go untuk Aggregator dan
Webhook, Rust untuk stream asinkron Feed, dan C\# untuk Notification. Compose
membangun image dari direktori masing-masing layanan, sehingga satu layanan
dapat dihentikan, dibangun ulang, dan dijalankan kembali tanpa menyentuh
layanan lain.

Setiap domain memiliki datanya sendiri (\emph{database per service}~\cite{dbperservice}):
Aggregator memiliki Canonical Store, Webhook Provider memiliki langganan dan
log pengiriman, Danger Notification memiliki notifikasi dan antrean push, dan
Event Live Feed memiliki buffer replay dalam memori. Tidak ada layanan yang
membaca tabel milik layanan lain. Domain penyampaian yang memerlukan data
hazard memperolehnya dari event, bukan dengan query ke Canonical Store.

\subsection{Canonical Store milik Aggregator}

Canonical Store adalah penyimpanan data hazard hasil normalisasi dan hanya
diakses oleh Aggregator. Implementasinya PostgreSQL 17 dengan skema
\emph{hybrid}: atribut yang sama untuk semua sumber disimpan sebagai kolom
bertipe, sedangkan atribut khusus sumber disimpan dalam satu kolom JSONB.

\begin{table}[htbp]
    \centering
    \small
    \caption{Struktur Canonical Store}
    \begin{tabularx}{\textwidth}{@{}P{0.25\textwidth}X@{}}
        \toprule
        \textbf{Tabel} & \textbf{Isi dan fungsi} \\
        \midrule
        \code{hazard_events} & Kolom bertipe \code{hazard_id}, \code{source}, \code{source_ref_id}, \code{hazard_type}, \code{severity}, \code{area_name}, koordinat, \code{occurred_at}, \code{ingested_at}; kolom \code{attributes} JSONB. Unique \code{(source, source_ref_id)} menjaga satu baris per kejadian sumber; indeks pada sumber, jenis, severity, dan \code{(occurred_at, hazard_id)} untuk filter dan paginasi. \\
        \code{ingest_watermarks} & Posisi polling, waktu sukses/percobaan terakhir, dan error per sumber; dasar status FRESH/STALE pada P2. \\
        \code{outbox_events} & Envelope event yang ditulis dalam transaksi yang sama dengan upsert hazard, menunggu dipublikasikan ke broker. \\
        \code{outbox_dead_letters} & Envelope yang gagal dipublikasikan setelah batas retry. \\
        \code{volcano_references} & Referensi nama dan koordinat gunung api milik BNPB untuk pemetaan PVMBG. \\
        \bottomrule
    \end{tabularx}
\end{table}

Ketika PVMBG menambahkan \code{confidence_level}, nilai tersebut masuk ke
\code{attributes} tanpa \code{ALTER TABLE}; record lama dan baru tersimpan pada
tabel yang sama dan dibaca dengan query yang sama. Filter lintas sumber tetap
memakai kolom bertipe dan indeks biasa. Skema diterapkan saat startup pada
\code{services/aggregator/internal/store/postgres.go}.

\section{Justifikasi dan Alternatif}

\subsection{Batas Layanan}

Alternatif pertama adalah \textbf{monolit modular}: semua domain dalam satu
proses dengan modul terpisah. Rancangan ini lebih sederhana dioperasikan,
tetapi kegagalan atau kebocoran memori pada satu modul (misalnya ribuan
koneksi feed) menjatuhkan seluruh fungsi, dan satu bahasa harus dipakai untuk
semua kebutuhan. Spesifikasi juga meminta satu layanan dapat dihentikan dan
di-deploy ulang tanpa mengganggu yang lain.

Alternatif kedua adalah \textbf{pemisahan lebih halus}, misalnya satu layanan
adapter per sumber (BMKG dan PVMBG) di depan Aggregator. Pemisahan ini
memberikan isolasi per sumber, tetapi pemetaan, upsert, dan penulisan outbox
perlu berada dalam satu transaksi agar tidak ada event yang hilang. Membagi
langkah tersebut ke beberapa layanan membutuhkan koordinasi transaksi
terdistribusi. Isolasi per sumber sudah dicapai dengan goroutine dan
watermark independen di dalam Aggregator (P2), sehingga pemisahan ini tidak
sebanding dengan biayanya.

\subsection{Penyimpanan Canonical Store}

\begin{table}[htbp]
    \centering
    \small
    \caption{Alternatif penyimpanan Canonical Store}
    \begin{tabularx}{\textwidth}{@{}P{0.21\textwidth}X X@{}}
        \toprule
        \textbf{Pilihan} & \textbf{Kelebihan} & \textbf{Biaya atau batasan} \\
        \midrule
        PostgreSQL relasional + JSONB (dipilih) & Atribut baru tidak memerlukan migrasi; kolom inti tetap bertipe dan terindeks; transaksi ACID mencakup hazard dan outbox; constraint unik menjaga identitas sumber. & Isi \code{attributes} tidak divalidasi skema database; query pada atribut dinamis memerlukan operator JSON atau indeks GIN tambahan. \\
        Relasional murni (kolom per atribut) & Semua field bertipe dan tervalidasi database. & Setiap field baru dari sumber memerlukan migrasi \code{ALTER TABLE} dan rilis ulang; bertentangan dengan kebutuhan menerima \code{confidence_level} saat runtime. \\
        MongoDB (dokumen) & Skema sepenuhnya fleksibel; payload sumber dapat disimpan apa adanya. & Field inti kehilangan jaminan tipe; constraint unik gabungan, transaksi multi-dokumen hazard--outbox, dan paginasi perlu dirancang ulang. Fleksibilitas yang dibutuhkan sudah dipenuhi JSONB. \\
        \bottomrule
    \end{tabularx}
\end{table}

PostgreSQL + JSONB dipilih karena kebutuhan Canonical Store terbagi dua:
bagian yang stabil dan sering difilter (sumber, jenis, severity, waktu) serta
bagian yang berubah mengikuti sumber. Skema hybrid melayani keduanya dalam satu
penyimpanan. Faktor penentu lainnya adalah transactional outbox pada P5, yang
membutuhkan upsert hazard dan penulisan event dalam satu transaksi atomik.

\section{Batasan Solusi}

Pemisahan image dan proses belum membuktikan bahwa setiap layanan dapat
direplikasi tanpa perubahan. Publisher outbox belum memakai mekanisme
claim/lock antar-worker pada query pengambilan record. Feed memiliki replay
dan epoch lokal per proses. SQLite webhook membatasi penulisan serta
memerlukan strategi berbagi state sebelum menjalankan beberapa replica.
Keterbatasan ini perlu diselesaikan sebelum menyatakan dukungan penskalaan
horizontal yang transparan.

Durability Canonical Store bergantung pada volume \code{aggregator_data}
dan prosedur backup; Compose belum menyediakan replikasi atau failover
PostgreSQL. Isi \code{attributes} tidak divalidasi oleh database, sehingga
konsistensi tipe atribut dinamis bergantung pada adapter Aggregator.

Penyimpanan hazard mempertahankan satu identitas kanonis per referensi
sumber. Pemisahan database dan penggunaan constraint tidak
dengan sendirinya menyediakan pemrosesan event tepat satu kali.

\section{Proof-of-Concept dan Keterlacakan}

Tangkapan layar pengujian layanan dan penyimpanan pada \code{docs/tex/evidence/P4/} mencakup rangkaian awal
dan tambahan \code{1-3-new.png} pada commit \code{4a1973b}. Ringkasan eksekusi otomatis pada
\code{docs/tex/evidence/problem_set_report.md}, dibuat 10 Oktober 2026,
mencatat tiga tes P4 lulus, tanpa failure maupun skipped. Cakupan kedua
jenis bukti dijelaskan berikut; label lulus tes tidak memperluas observasi
melampaui assertion dan screenshot yang tersedia.

\subsection{Pemisahan Proses dan Pengoperasian Layanan}

Pada pengujian pemisahan proses, operator menghentikan \code{notification-service}. Halaman
notifikasi kemudian menampilkan respons 503 untuk endpoint yang bergantung
pada layanan tersebut, sementara feed tetap aktif dan status sumber tetap
tersedia. Setelah layanan dibangun dan dijalankan kembali, halaman
notifikasi kembali menampilkan data. Urutan ini direkam pada
\code{P4/1-1.png} sampai \code{P4/1-4.png}.

\evidencefigure{evidence/P4/1-1.png}
{Penghentian container Danger Notification secara terpisah (\code{P4/1-1.png}).}
{fig:p4-stop-notification}

Grafik health memperlihatkan interval notification tidak tersedia,
sedangkan Aggregator, feed, gateway, dan database tetap terpantau aktif.
Webhook memang sedang tidak aktif pada interval grafik tersebut; gambar
ini tidak digunakan untuk menyatakan semua layanan hidup bersamaan.

\evidencefigure{evidence/P4/1-5.png}
{Observasi health ketika notification dihentikan dan dipulihkan (\code{P4/1-5.png}).}
{fig:p4-health-independence}

Tes otomatis \code{tests/test_problem_4.py} melengkapi observasi dengan
menghentikan satu layanan, meminta hazard BMKG melalui API, memeriksa
respons 200 beserta data, lalu menjalankan kembali layanan. Ini mendukung
ketersediaan jalur baca independen selama penghentian yang diuji.
Bukti tambahan \code{P4/1-3-new.png} memakai perintah berikut:
\begin{lstlisting}
docker compose --env-file .env --env-file .env.auth up --build --no-deps notification-service
\end{lstlisting}
Output memperlihatkan 24 langkah build hanya untuk Notification, image
berstatus Built, dan container Notification berstatus Created lalu di-attach.
Sebagian besar langkah memakai build cache. Opsi \code{--no-deps} membatasi
operasi pada target sehingga dependency tidak ikut dibangun atau dijalankan
oleh perintah ini. Bukti ini menggantikan \code{P4/1-3.png} sebagai rujukan
build satu layanan; screenshot lama tetap tersimpan sebagai riwayat.

\evidencefigure{evidence/P4/1-3-new.png}
{Build dan pembuatan container Notification dengan \code{--no-deps}; output hanya menampilkan layanan target (\code{P4/1-3-new.png}).}
{fig:p4-rebuild-no-deps}

Grafik health dan pembacaan API pada rangkaian awal melengkapi observasi
ketersediaan jalur independen. Screenshot build baru tidak memuat
perbandingan ID/start time seluruh container atau health setelah attach,
sehingga tidak digunakan sebagai catatan lengkap riwayat restart setiap
container di luar interval yang direkam.

\subsection{Atribut Dinamis dan Koeksistensi Skema}

Setelah kontrol PVMBG diubah ke versi 2 dan ditunggu 60 detik, pembacaan
100 hazard PVMBG melalui Canonical Store menemukan 25 record tanpa
\code{confidence_level} dan 75 record dengan atribut tersebut. Screenshot
menampilkan contoh kedua bentuk atribut, termasuk nilai confidence 0,69.
Pemeriksaan tabel memperlihatkan bahwa \code{attributes} bertipe JSONB;
\code{confidence_level} tidak menjadi kolom terpisah.

\evidencefigure{evidence/P4/2-1.png}
{Record PVMBG sebelum dan sesudah penambahan atribut tersedia bersamaan melalui API (\code{P4/2-1.png}).}
{fig:p4-schema-coexistence}

\evidencefigure{evidence/P4/2-2.png}
{Struktur tabel hazard menggunakan JSONB untuk atribut sumber (\code{P4/2-2.png}).}
{fig:p4-jsonb-schema}

Hasil ini mendukung rancangan penyimpanan fleksibel dan koeksistensi data
v1/v2. Tes otomatis memeriksa keberadaan kedua kelompok record serta tipe
numerik confidence. Screenshot DDL yang tersedia adalah satu snapshot
struktur, bukan pasangan hasil sebelum/sesudah. Oleh sebab itu, tidak ada
angka perubahan DDL atau hasil perbandingan hash schema yang diklaim.

\subsection{Kepemilikan Data dan Batas Pengujian}

Tangkapan layar \code{P4/3.png} memperlihatkan inspeksi konfigurasi serta keberhasilan
pembacaan status sumber melalui API. Pada kode aplikasi, operasi Canonical
Store ditangani Aggregator; consumer memperoleh data lewat event atau API.
Pemeriksaan otomatis mencari variabel \code{CANONICAL_STORE_DSN} secara
literal dalam blok Compose layanan lain. Pemeriksaan ini menunjukkan
konvensi konfigurasi yang diperiksa, tetapi tidak membuktikan pembatasan
akses database pada tingkat jaringan atau hak pengguna.

Compose tidak lagi menyuntikkan seluruh \code{.env} ke container; setiap
layanan menerima variabel yang diperlukannya saja, dan DSN Canonical Store
hanya diberikan kepada Aggregator. Namun, database berada pada jaringan bersama
dan port PostgreSQL diekspos ke host; collector observabilitas mempunyai
koneksi PostgreSQL untuk metrik. Dengan demikian, istilah kepemilikan
pada laporan ini berarti pemisahan tanggung jawab aplikasi. Pembatasan
credential per layanan, role metrik khusus, dan jaringan database perlu
diterapkan bila kriteria akses eksklusif hendak dijamin secara teknis.

\begin{table}[htbp]
    \centering
    \small
    \caption{Cakupan bukti P4 dan batas interpretasi}
    \begin{tabularx}{\textwidth}{@{}P{0.25\textwidth}X@{}}
        \toprule
        \textbf{Objek} & \textbf{Hasil dan batas bukti} \\
        \midrule
        Independensi proses & Stop dan pemulihan diamati; bukti baru memperlihatkan build target dengan \code{--no-deps}. Health jalur lain tersedia pada rangkaian awal. Identitas semua container belum dibandingkan. \\
        Atribut dinamis & 25 record tanpa confidence dan 75 dengan confidence terbaca; kolom JSONB diperlihatkan. \\
        Kepemilikan Canonical Store & Jalur aplikasi melalui Aggregator/event ditunjukkan; konfigurasi belum menjamin isolasi credential dan jaringan. \\
        Persistensi lintas restart & Volume PostgreSQL, webhook, dan notification tersedia pada kode. Tangkapan layar pengujian layanan dan penyimpanan tidak membandingkan ID/jumlah data ketiganya sebelum dan sesudah restart. \\
        \bottomrule
    \end{tabularx}
\end{table}

Hasil historis pengujian container mock pada
\code{docs/tex/evidence/mock-container-2026-09-30.json} tetap hanya berlaku untuk persistensi database mock yang
ketika itu diuji; tidak digunakan sebagai pengganti pengujian volume
Canonical Store, webhook, atau notification. Pengujian unit Aggregator
pada \code{docs/tex/evidence/aggregator-unit-2026-10-09.jsonl}
mendukung transformasi atribut, sedangkan ringkasan eksekusi otomatis dan
tangkapan layar pengujian layanan dan penyimpanan menambahkan pengamatan
pada stack berjalan. Sumber bukti terdaftar pada
\code{docs/tex/evidence/manifest.json}; petunjuk pengujian tersedia pada
\code{docs/e2e-testing.md}.
